% =============================================================
%  Symbols
%  Collects the symbols used across the report, in the notation of
%  Section 3 (sending/receiving end, primed quantities per unit
%  length).  Matching slide: presentation/design/ac-transmission/
%  symbols.tex.
% =============================================================
\section{Symbols}
\label{sec:symbols}

Quantities per unit length carry a prime, $R'$, $L'$, $G'$, $C'$; the same
letter without the prime is the quantity of the whole line. Complex
quantities are underlined. Terminal quantities take the subscript $S$ at the
sending end and $R$ at the receiving end, and voltages are line-to-line
unless stated otherwise.

\begin{tabularx}{\linewidth}{@{}l X@{}}
  \toprule
  \multicolumn{1}{c}{\textbf{Symbol}} & \textbf{Description} \\
  \midrule
  $\underline{A}, \underline{B}, \underline{C}, \underline{D}$
          & Two-port (ABCD) parameters of the line \\
  $C'$    & Capacitance per unit length \\
  $d$     & Conductor diameter \\
  $D$     & Phase spacing of the tower head \\
  $f$     & System frequency; $f_{er}$ the electrical resonance frequency \\
  $G'$    & Conductance per unit length \\
  $GMD$   & Geometric mean distance between the phases, $D_m$ \\
  $GMR$   & Geometric mean radius of a sub-conductor; $GMR_b$ of the bundle \\
  $H$     & Inertia constant of a synchronous machine \\
  $I$     & Phase current; $\underline{I}_S$, $\underline{I}_R$ at the terminals \\
  $k$     & Degree of series compensation, $X_C = k \cdot X_L$ \\
  $K_s$   & Synchronising power coefficient, $\partial P / \partial \delta$ \\
  $l$     & Line length; $x$ the distance from the sending end \\
  $L'$    & Inductance per unit length \\
  $m$     & Degree of shunt compensation \\
  $n$     & Turns ratio of an ideal transformer; number of sections in a cascade \\
  $P$     & Active power; $P_{nat}$ natural power, $P_{max}$ maximum transfer \\
  $Q$     & Reactive power \\
  $r_b$   & Equivalent radius of the bundle for the capacitance \\
  $R'$    & Resistance per unit length \\
  $s$     & Bundle spacing \\
  $S_{sc}$ & Short-circuit level at a terminal \\
  $U$     & Line-to-line voltage; $\underline{U}_S$, $\underline{U}_R$ at the
            terminals, $U_n$ nominal \\
  $X$     & Series reactance of the path; $X'$ per unit length, $X_C$ of the
            series bank, $X_{tot}$ of the whole path between the internal
            voltages \\
  $\underline{Y}'$ & Shunt admittance per unit length, $G' + j \cdot \omega \cdot C'$ \\
  $\underline{Z}'$ & Series impedance per unit length, $R' + j \cdot \omega \cdot L'$ \\
  $\underline{T}$ & Chain (ABCD) matrix of a two-port \\
  $\underline{Z}_{eq}$ & Series branch of the equivalent $\pi$;
            $\underline{Y}_{eq}$ its shunt branch \\
  $\underline{Z}_c$ & Image (characteristic) impedance of a two-port,
            $\sqrt{\underline{B}/\underline{C}}$ \\
  $\underline{Z}_{in}$ & Impedance seen at port~1 of a terminated two-port \\
  $\underline{Z}_L$ & Impedance terminating a two-port \\
  $\underline{Z}_w$ & Wave (surge) impedance, $\sqrt{\underline{Z}'/\underline{Y}'}$ \\
  \midrule
  $\alpha$   & Attenuation constant \\
  $\beta$    & Phase constant; $\beta \cdot l$ the electrical length \\
  $\underline{\gamma}$ & Propagation constant, $\alpha + j \cdot \beta$ \\
  $\delta$   & Transmission angle between the terminal voltages \\
  $\varepsilon_0$ & Permittivity of free space \\
  $\mu_0$    & Permeability of free space \\
  $\underline{\theta}$ & Image transfer constant of a two-port;
            $\underline{\gamma} \cdot l$ for a line \\
  $\varphi$  & Power-factor angle \\
  $\omega$   & Angular frequency, $2 \cdot \pi \cdot f$ \\
  \bottomrule
\end{tabularx}

\FloatBarrier
